\documentclass[sigconf,anonymous,review]{acmart}

\setcopyright{none}
\settopmatter{printfolios=true}
\renewcommand\footnotetextcopyrightpermission[1]{}
\pagestyle{plain}
\acmConference[DAC '26]{63rd ACM/IEEE Design Automation Conference}{July 26--29, 2026}{Long Beach, CA, USA}
\acmYear{2026}
\copyrightyear{2026}
\acmISBN{}
\acmDOI{}

\newcommand{\system}{\textsc{PersistRoute}}
\newcommand{\rrr}{rip-up-and-reroute}
\newcommand{\todoresult}[1]{\textbf{[BOARD RESULT REQUIRED: #1]}}
\setlength{\emergencystretch}{1.5em}

\begin{document}

\title{PersistRoute: Transactional Persistent-State FPGA Routing}

\author{Anonymous Author(s)}
\affiliation{\institution{Anonymous Institution}}

\begin{abstract}
Negotiated-congestion FPGA routing repeatedly traverses a large routing-resource
(RR) graph while updating occupancy, history cost, route trees, and timing
criticality. Accelerating isolated searches leaves state movement and global
synchronization on the critical path. \system{} instead keeps production
routing state resident in HBM across complete rip-up-and-reroute iterations.
Four Search compute units speculatively generate versioned candidate paths with
banked MultiQueues; one deterministic Commit unit validates and applies those
transactions in a fixed order. A host control layer preserves PathFinder
semantics and invokes VPR/Tatum after each iteration to refresh per-connection
setup criticality. The complete 32-HBM-port design routes and closes timing on
an AMD Alveo U280 at 180\,MHz, using 16.48\% of LUTs. Hardware emulation on a
real VPR RR-graph subset is byte-identical to the HBM C-model for concurrent
search and two-iteration negotiated congestion. Ten paired VTR, Titan23, and
Koios cases further establish a deterministic multicore baseline with
1.356$\times$ geometric-mean routing-stage scaling. These results validate the
architecture and its routing semantics; physical-board end-to-end performance
and energy remain required for an acceleration claim.
\end{abstract}

\begin{CCSXML}
<ccs2012>
 <concept>
  <concept_id>10010520.10010575.10010755</concept_id>
  <concept_desc>Computer systems organization~Reconfigurable computing</concept_desc>
  <concept_significance>500</concept_significance>
 </concept>
 <concept>
  <concept_id>10002950.10003624.10003633</concept_id>
  <concept_desc>Mathematics of computing~Graph algorithms</concept_desc>
  <concept_significance>300</concept_significance>
 </concept>
</ccs2012>
\end{CCSXML}
\ccsdesc[500]{Computer systems organization~Reconfigurable computing}
\ccsdesc[300]{Mathematics of computing~Graph algorithms}
\keywords{FPGA routing, PathFinder, HBM, speculative execution, deterministic commit}

\maketitle

\section{Introduction}

Routing dominates difficult FPGA compilations because each signal connection
performs an irregular shortest-path search over a large RR graph and because
negotiated congestion revisits the same state for many iterations. PathFinder
turns illegal sharing into progressively higher present and historical costs
\cite{mcmurchie1995pathfinder}; modern VPR adds connection-level routing,
lookahead, route trees, and timing-driven criticalities \cite{vtr8}. These
features improve quality of results (QoR), but they also expose an unfavorable
machine behavior: latency-sensitive priority queues generate random accesses,
connections overlap in space, and every iteration synchronizes on shared
occupancy and timing state.

Prior work attacks complementary pieces of this problem. Deterministic and
partitioned CPU routers exploit inter-net or inter-connection parallelism
\cite{moctar2018galois,zhou2020connection}; MultiQueue routing relaxes strict
priority ordering inside one search \cite{singer2024multiqueue}; and recent
commercial routers allow or manage overlap among concurrently routed nets
\cite{zang2024cufr,zang2024potter}. These advances do not by themselves define
how a hardware accelerator should retain the authoritative RR state across
iterations, how speculative paths should become deterministic updates, or how
real timing analysis should feed the next hardware round.

\system{} addresses that systems gap. It keeps the graph, resource state,
queues, work arrays, paths, and transaction logs in U280 HBM. Four Search
compute units (CUs) read versioned resource snapshots and emit candidate paths
without mutating authoritative state. A single Commit CU checks the recorded
versions and applies valid transactions in a fixed order; rejected work enters
a replay queue. This separation admits overlapping connection searches without
races while preserving a reproducible state transition. The same 32-byte task
ABI drives a C++ reference, an HBM C-model, and an XRT backend.

This paper makes four contributions:
\begin{itemize}
 \item a cross-iteration HBM state organization for production-scale RR graphs,
 including occupancy, congestion cost, route trees, queues, and version state;
 \item a two-level execution scheme combining concurrent connection searches
 with banked MultiQueues inside each search;
 \item a deterministic transaction protocol that separates speculative search
 from fixed-order version-checked commit and replay; and
 \item an auditable VPR/FPGAIF-to-XRT flow with per-iteration VPR/Tatum timing
 feedback and layered C-model, RTL-emulation, implementation, and legality
 validation.
\end{itemize}

We deliberately separate validated facts from unavailable measurements. The
full U280 design is implemented and timing-closed, and RTL emulation matches the
software model, but this VMware environment has no physical U280. Consequently,
we do not report emulation wall time as speedup or claim energy improvement.

\section{Background and Motivation}

\subsection{Negotiated-congestion routing}

Let $G=(V,E)$ be the directed RR graph. A connection task $k$ has source
$s_k$, sink $t_k$, timing criticality $\alpha_k$, and a bounding box. Search
minimizes a timing/congestion objective of the form
\begin{equation}
 f_k(v)=g_k(v)+(1-\alpha_k)h_{\mathrm{cong}}(v)+
             \alpha_k h_{\mathrm{delay}}(v),
\end{equation}
where $g_k$ is accumulated path cost and the two lookahead components estimate
congestion and delay. A resource $v$ with occupancy $o_v$ and capacity $c_v$
is overused when $o_v>c_v$. After an iteration, PathFinder increases present
cost for current overuse and historical cost for persistent overuse. Nets
touching illegal resources are ripped up and searched again. A correct
accelerator must therefore support more than independent point-to-point maze
queries: it must preserve route trees, update shared negotiated costs, and
repeat until all connections are routed with zero overuse.

Timing-driven routing introduces a second feedback loop. Connection delay
changes after each route. Static timing analysis (STA) then updates slack and
criticality, changing priorities and edge costs in the next round. Holding
criticality fixed is convenient for a kernel demonstration but does not match
a production timing-driven router.

\subsection{Why conventional offload is insufficient}

Three observations motivate persistent transactional execution.
\textit{First}, RR graphs are much larger than the set of paths changed in one
iteration. Re-transferring graph and occupancy state for every batch can erase
search acceleration. \textit{Second}, spatial partitioning alone loses
parallelism on high-fanout or long connections whose bounding boxes overlap.
Allowing overlap shifts the problem from scheduling to conflict resolution.
\textit{Third}, mutating occupancy as each CU finishes makes output dependent on
memory arbitration and completion time. Reproducibility is important both for
CAD debugging and for paired QoR evaluation.

\system{} therefore treats Search as speculative and Commit as authoritative.
The hardware retains state across all rounds; only a compact timing snapshot
crosses the host boundary when real STA is requested. This structure targets
the random-access and synchronization bottlenecks without replacing the FPGA
routing problem by an ASIC grid-routing abstraction.

\section{PersistRoute Architecture}

\begin{figure*}[t]
\centering
\setlength{\tabcolsep}{3pt}
\begin{tabular}{c@{\,$\rightarrow$\,}c@{\,$\rightarrow$\,}c@{\,$\rightarrow$\,}c@{\,$\rightarrow$\,}c}
\fbox{\shortstack{VPR / FPGAIF\\converter}} &
\fbox{\shortstack{Host controller\\epochs, timing, replay}} &
\fbox{\shortstack{4$\times$ Search CU\\8-bank MultiQueue}} &
\fbox{\shortstack{Transaction HBM\\paths + read versions}} &
\fbox{\shortstack{1$\times$ Commit CU\\fixed-order validation}}
\\[5pt]
\multicolumn{5}{c}{\fbox{\shortstack{32 HBM pseudo-channels: RR nodes/edges, resources, work arrays, queues, route trees, and logs}}}
\end{tabular}
\caption{\system{} separates speculative path discovery from authoritative,
deterministic state mutation while retaining routing state across complete RRR
iterations.}
\Description{A left-to-right architecture diagram connects the VPR or FPGA
Interchange converter, host epoch controller, four parallel Search compute
units, transaction storage, and one deterministic Commit compute unit above
32 HBM pseudo-channels.}
\label{fig:architecture}
\end{figure*}

\subsection{Production data representation}

The front end converts VPR RR-graph XML or FPGA Interchange device resources
into structure-of-arrays binaries. Node and edge indices are 32 bits, avoiding
the small fixed graph limits common in on-chip demonstrations. A task is a
fixed 32-byte record containing connection identity, source, sink, net,
criticality, bounding box, and control flags. Keeping this ABI unchanged across
the CPU reference, HBM model, and XRT backend makes binary inputs and outputs
directly comparable.

The graph is immutable during routing. Mutable resource records hold capacity,
occupancy, present/history cost, owner fast-path state, and a monotonically
increasing version. Work arrays and queues live in HBM rather than BRAM, so
their size follows graph scale. The U280 platform exposes 32 HBM
pseudo-channels; the linked design connects four Search CUs and one Commit CU
through 32 kernel master ports. Read-mostly graph banks are separated from
write-heavy transaction and resource banks to reduce interference.

\subsection{Banked MultiQueue Search}

Each Search CU runs an A*-style connection search. A strict global heap
serializes push/pop operations and creates a hard random-access dependence.
Instead, an eight-bank MultiQueue hashes candidates into independent binary
heaps. A pop samples bank heads and removes a low-cost candidate; relaxed
ordering increases available memory-level and pipeline parallelism while the
best-known cost array suppresses dominated expansions. Deterministic bank
selection and task-local tie breaking make repeated C-model executions
reproducible even though the ordering is relaxed.

Search reads resource occupancy, costs, and versions but does not claim
resources. Its output is a candidate node path plus a compact read/write set.
For each resource on the candidate path, the transaction records the observed
version. Logging the entire explored region is unnecessary for commit
feasibility and, on large graphs, can exceed the 16-bit per-transaction count;
changes off the chosen path may affect optimality but cannot invalidate that
path's capacity check. Failed
searches also produce explicit status, preventing missing paths from being
mistaken for empty legal routes. Because candidates are isolated, four CUs can
search overlapping regions in the same wave.

\subsection{Versioned deterministic commit}

The host orders completed candidates by descending timing criticality, then
$(\text{priority},\text{task-id})$, and
streams them to one Commit CU. For transaction $T_k$ with read set $R_k$, commit
succeeds iff
\begin{equation}
 \forall v\in R_k:\quad \mathrm{version}_{\mathrm{now}}(v)
                  =\mathrm{version}_{T_k}(v).
\end{equation}
On success, Commit applies route-tree and occupancy changes, updates ownership,
and increments the affected versions. A mismatch rejects the whole transaction
without a partial update; the task is replayed against the new state. Fixed
ordering converts nondeterministic Search completion into a deterministic
sequence of authoritative mutations.

At the end of an iteration, the controller identifies overused resources,
rips up affected connections, updates negotiated costs, and forms the next
epoch. Termination requires every connection to have a path and every resource
to satisfy capacity. The architecture thus accelerates the repeated control
loop, not only a one-shot path-search kernel.

\subsection{Protocol invariants and failure handling}

The transaction split is useful only if it preserves serial state semantics.
We enforce four invariants. (I1) Search never writes authoritative occupancy,
owner, congestion, or route-tree state. (I2) Every resource read by a candidate
is paired with the version observed by Search. (I3) Commit either applies all
of a transaction's resource and route-tree updates or applies none. (I4)
Successful transactions are linearized in the published deterministic order.
I1 removes write races among Search CUs. I2 and I3 detect interference and
prevent partially installed paths. I4 defines one reproducible serial history
for an epoch, independent of CU latency.

\begin{figure}[t]
\small
\begin{tabular}{@{}rl@{}}
1: & order by criticality, priority, and connection ID \\
2: & while an ordered wave remains: \\
3: & \quad launch up to four isolated searches \\
4: & \quad collect candidate paths and read versions \\
5: & \quad for each candidate in fixed order: \\
6: & \qquad if every version matches, commit atomically \\
7: & \qquad otherwise enqueue the task for replay \\
8: & update overuse and negotiated congestion costs \\
9: & rip up illegal connections; begin the next epoch \\
10: & stop when all tasks route and overuse is zero
\end{tabular}
\caption{Deterministic speculative routing control. ``Atomically'' denotes one
logical transaction; Commit may stream its validated writes.}
\Description{Ten-line pseudocode orders connections, launches up to four
isolated searches, validates candidates in a fixed order, replays conflicts,
updates congestion costs, and stops at a legal zero-overuse routing.}
\label{fig:protocol}
\end{figure}

Resource-version overflow is made explicit by a sufficiently wide monotonic
counter in the production record; a run must fail rather than silently wrap.
Malformed paths, out-of-range node IDs, missing task results, queue exhaustion,
and path-length overflow produce errors at the ABI boundary. A search that
cannot reach its sink produces a failed transaction and remains eligible for a
later epoch after congestion costs change. The CPU and accelerated validators
independently recheck edge adjacency, path endpoints, occupancy, route-tree
membership, and zero final overuse.

Determinism does not imply that relaxed MultiQueue expansion visits exactly the
same nodes as a strict heap. It means that identical inputs and configuration
produce identical \system{} outputs, and speculative completion order cannot
select a different authoritative state. QoR equivalence to VPR is a separate
empirical requirement, measured by legality, wirelength, CPD, and iteration
count rather than inferred from the protocol.

\section{Timing and Flow Integration}

\subsection{Per-iteration real timing feedback}

After every complete iteration, the host reconstructs each connection's RR
node path and sums quantized edge delays. Edge delay uses VPR's buffered local
$T_{\mathrm{linear}}$ RC expression; pass switches use a documented local RC
approximation. The host writes one delay record per connection and invokes a
patched VPR analysis run. The patch maps $(\text{net},s,t)$ back to the packed
netlist, overwrites the corresponding entries of VPR's \texttt{net\_delay}
matrix, and runs the normal \texttt{AnalysisDelayCalculator} and Tatum setup
analysis. Returned slack, setup criticality, and critical-path delay are
validated one-to-one against task IDs and quantized into the next iteration's
unchanged 32-byte task record.

Fail-closed validation rejects missing, duplicated, malformed, or ambiguous
connections. Unrouted connections retain the delay of VPR's legal reference
route so a failed intermediate search does not permanently eliminate a task.
The current oracle supports clustered routing and intentionally rejects flat
routing. Although host STA is not persistent hardware state, its periodic,
audited interface preserves a real timing graph rather than substituting a
synthetic criticality schedule.

\subsection{Scheduling and memory traffic}

The host groups ready tasks into four-CU waves while preserving a global commit
order. Search load imbalance is expected: path length and explored volume vary
by orders of magnitude even among connections of the same net. Static spatial
partitioning would leave CUs idle behind a long connection. Wave scheduling
instead assigns the next ready task to the next free Search CU; determinism is
recovered at Commit rather than by synchronizing expansion. The accelerated
iteration log records launched, committed, failed, and replayed task counts,
overuse, and state hashes, permitting post-run utilization and conflict
analysis.

\begin{table}[t]
\centering
\small
\begin{tabular}{lll}
\toprule
State & Access pattern & Lifetime \\
\midrule
RR nodes/edges & read-mostly & all iterations \\
Resources & Search read, Commit write & all iterations \\
Best-cost/parent & Search read/write & one task \\
Queue banks & Search read/write & one task \\
Candidate path & Search write, Commit read & one attempt \\
Route tree & Commit/rip-up read & all iterations \\
Timing snapshot & host/VPR exchange & one iteration \\
\bottomrule
\end{tabular}
\caption{State access and lifetime. Persistent entries avoid graph-scale
reinitialization between RRR iterations.}
\label{tab:lifetime}
\end{table}

Table~\ref{tab:lifetime} separates persistent and task-local traffic. If the RR
graph contains $N$ nodes and $M$ edges but an iteration touches $L$ path nodes,
a full-state offload moves $\Theta(N+M)$ immutable data per batch. \system{}
transfers that graph once and thereafter communicates task descriptors,
transaction/path records proportional to explored work, and periodic timing
records proportional to the connection count. This does not eliminate HBM
random access; it changes repeated host/device movement into on-device bank
traffic. A physical-board ablation must quantify whether that reduction
dominates PCIe setup and host STA.

\subsection{End-to-end and commercial interfaces}

The academic path accepts VPR architecture, netlist, placement, route, and RR
graph artifacts. The commercial path uses RapidWright/FPGA Interchange to
extract and replay routed programmable interconnect points. The same converter
tests enforce node/edge/task record sizes and hashes. On the FPGA24
\texttt{logicnets\_jscl} design, the bridge performs a non-empty six-PIP
round-trip and Vivado routes all 27,961 routable nets with zero routing errors
and zero DRC violations. Because that public input contains no user timing
constraints, this result establishes legality, not timing closure. RWRoute
\cite{zhou2021rwroute}, CUFR \cite{zang2024cufr}, and Vivado provide
same-input commercial references.

\section{Evaluation}

\subsection{Methodology}

We use VTR 9, Titan 2.0/Titan23 \cite{murray2013titan,talaei2023titan2}, and
Koios 2.0 \cite{arora2021koios}. Each paired CPU baseline invokes the same
TBB-enabled VPR executable; only \texttt{--num\_workers} changes from 1 to 16
(on an eight-logical-CPU VM). Architecture, synthesized netlist, SDC when
available, seed, channel width, and router options are identical. A pair is
accepted only if both runs pass and route SHA-256, wirelength, critical-path
delay (CPD), and iteration count match. We report both full VPR wall time and
routing-stage time to expose Amdahl's-law limits.

Hardware synthesis, co-simulation, link, placement, and routing use Vitis,
Vitis HLS, and Vivado 2022.2 for an AMD XCU280 device. The linked
system has four Search and one Commit CU. Hardware-emulation outputs are
compared byte-for-byte with the HBM C-model. Since no U280 is passed through to
the VM, emulation time is excluded from speedup and energy results.

\begin{table*}[t]
\centering
\small
\begin{tabular}{llrrrrrr}
\toprule
Suite & Circuit & Blocks & Nets & Route iters. & CPD (ns) & Flow speedup & Route speedup \\
\midrule
VTR & diffeq & 1974 & 1935 & 13 & 6.598 & 1.025$\times$ & 1.333$\times$ \\
VTR & frisc & 4561 & 4445 & 12 & 13.588 & 1.022$\times$ & 1.318$\times$ \\
VTR & elliptic & 4969 & 4855 & 12 & 8.792 & 1.037$\times$ & 1.421$\times$ \\
VTR & clma & 8542 & 8460 & 15 & 11.608 & 1.055$\times$ & 1.421$\times$ \\
Titan23 & neuron & 86875 & 119888 & 15 & 7.603 & 1.090$\times$ & 1.293$\times$ \\
Titan23 & stereo vision & 94090 & 127090 & 15 & 7.517 & 1.088$\times$ & 1.417$\times$ \\
Koios 2.0 & conv layer & 28067 & 33180 & 14 & 5.128 & 1.089$\times$ & 1.420$\times$ \\
Koios 2.0 & attention layer & 36457 & 45949 & 15 & 8.823 & 1.017$\times$ & 1.262$\times$ \\
Koios 2.0 & SpMV & 14407 & 16311 & 13 & 6.153 & 0.996$\times$ & 1.297$\times$ \\
Koios 2.0 & softmax & 11781 & 12958 & 18 & 9.190 & 1.049$\times$ & 1.392$\times$ \\
\midrule
\multicolumn{6}{r}{Geometric mean} & 1.046$\times$ & 1.356$\times$ \\
\bottomrule
\end{tabular}

\caption{Deterministic paired VPR CPU baseline. ``Flow'' includes packing and
placement; ``Route'' isolates routing. Every pair has identical route hash,
wirelength, CPD, and routing-iteration count.}
\label{tab:cpu}
\end{table*}

\subsection{CPU baseline characterization}

Table~\ref{tab:cpu} covers four VTR circuits, both Titan23 designs available in
the local Titan release, and four structurally different Koios workloads. The
ten pairs are bitwise route-identical. Routing-stage geometric-mean scaling is
1.356$\times$, whereas whole-flow scaling is only 1.046$\times$ because packing,
placement, and serial analysis dominate. The largest route-stage gain is
1.421$\times$; \texttt{spmv} slightly slows down end-to-end (0.996$\times$)
despite a 1.297$\times$ routing-stage gain. This negative case is retained.
The result establishes a reproducible strong CPU reference and shows why an
accelerator paper must report total time rather than a selected kernel alone.

\subsection{Implementation results}

\begin{table}[t]
\centering
\small
\begin{tabular}{lrrr}
\toprule
Resource & Used & Available & Utilization \\
\midrule
LUT & 186821 & 1133280 & 16.48\% \\
Register & 276771 & 2266560 & 12.21\% \\
BRAM & 108 & 1776 & 6.08\% \\
DSP & 52 & 8304 & 0.63\% \\
\bottomrule
\end{tabular}

\caption{Post-synthesis dynamic-region resources for the complete four-Search,
one-Commit U280 system.}
\label{tab:resources}
\end{table}

The complete batched, four-replica hardware flow passes system link, block
design generation, synthesis, optimization, placement, routing, bitstream
generation, and XCLBIN packaging. At the requested 180\,MHz data clock, routed setup WNS is
+0.016\,ns with TNS 0; hold WHS is +0.006\,ns with THS 0. There are no setup or
hold failing endpoints. Table~\ref{tab:resources} shows 16.48\% LUT and 12.21\%
register use, leaving capacity for additional Search CUs or larger local
caches, subject to HBM-port and timing constraints. Search is the larger
kernel; the final system timing result, rather than an older out-of-context
estimate, is the implementation claim.

\begin{table}[t]
\centering
\small
\begin{tabular}{p{0.30\columnwidth}p{0.61\columnwidth}}
\toprule
Test & Observed result \\
\midrule
HLS C/RTL & Five tops pass C simulation and XSIM co-simulation. \\
Four-CU wave & Four real-VPR connections launch together; legal in one round;
all five output classes match the C-model. \\
Conflict/replay & Two connections create one overused resource in round 1 and
reach zero overuse in round 2; outputs match byte-for-byte. \\
Timing loop & Full \texttt{diffeq} and \texttt{frisc} C-model routes use
per-round VPR/Tatum STA; both close legally. \\
VPR route check & Exported \texttt{frisc}/\texttt{elliptic} routes pass native
routing consistency and feasibility checks. \\
Commercial check & 27,961/27,961 routable nets; zero routing errors and DRC
violations after Vivado route. \\
\bottomrule
\end{tabular}
\caption{Layered functional and flow validation.}
\label{tab:validation}
\end{table}

\subsection{Functional closure and timing feedback}

In the four-task RTL case, all four Search CUs launch in one wave and Commit
produces a legal one-round result. Paths, route trees, ordinary iteration logs,
accelerated iteration logs, and final resource bytes match the HBM C-model.
The conflict case starts with one overused resource after the first commit,
then exercises host rip-up, reroute, and replay; the second iteration reaches
zero overuse and again matches all output classes. Simulator processes exit
normally with no XSIM error or fatal messages.

The full \texttt{diffeq} timing experiment has 32,920 RR nodes, 344,820 edges,
and 1,973 connections. CPU routing legally converges in three iterations; the
four-PE HBM C-model converges in twelve. VPR/Tatum successfully analyzes every
round, returning final CPDs of 7.164 and 7.274\,ns, respectively. This proves
that changing connection paths affect the real timing graph and the next
round's priorities.

On the 74,928-node, 5,767-connection \texttt{frisc} graph, a fixed task order
with base cost 256 finishes 5.73\% above VPR's 13.588-ns CPD. Normalizing the
base cost to 64 reduces the gap to 3.45\%. One frozen policy---base cost 64,
three timing-aware reroute rounds, and deterministic criticality ordering---
yields feedback-oracle CPDs of 13.697\,ns (+0.803\%) for \texttt{frisc} and
8.836\,ns (+0.499\%) for \texttt{elliptic}. Deterministic route export then
passes VPR's native consistency and feasibility checks. Native CPDs are
13.708\,ns (+0.884\%) and 8.840\,ns (+0.549\%); wirelengths are 34,395
(+3.656\%) and 28,385 (+3.864\%). Both paired C-model runs are zero-overuse
and byte-identical. Thus two designs meet the CPD gate but not the
pre-registered 2\% wirelength gate; this is not a ten-design QoR result.
A preliminary static RR channel-length penalty worsens \texttt{frisc}
wirelength, so we reject scalar edge reweighting as the cleanup mechanism.
Instead, a post-legal cleanup snapshots and reroutes one net at a time,
committing only legal tree-wirelength reductions and rolling back all other
candidates. Two passes yield native wirelength/CPD deltas of
+3.134\%/+0.876\% on \texttt{frisc} and +2.646\%/+0.084\% on
\texttt{elliptic}. Geometric-mean wirelength degradation falls from 3.760\%
to 2.890\%, an improvement that still misses the 2\% gate.

\subsection{Scale coverage and stress behavior}

The converted real graphs exercise 32-bit production indices rather than only
legacy on-chip limits. \texttt{diffeq} produces 32,920 nodes, 344,820 edges,
and 1,973 connection tasks. The larger \texttt{clma} conversion contains
155,280 nodes, 1,803,328 edges, and 13,458 tasks. On the CPU golden model,
\texttt{clma} routes every connection and closes legally in three iterations,
performing 23,060 searches, 224.36 million expansions, and 1.479 billion edge
visits. These counts illustrate both the opportunity for parallel expansion
and the external-memory pressure hidden by tiny kernel tests.

The academic matrix spans 1,974--94,090 blocks and 1,935--127,090 nets.
Titan23 is largest by netlist size, while Koios exercises heterogeneous
hard-memory/DSP architecture models and markedly different front-end costs.
The 3.0-GiB peak RSS of \texttt{attention\_layer} and its small whole-flow
scaling show that routing acceleration must coexist with large serial
packing/placement components. We therefore treat CPU stage scaling as workload
characterization, not as a proxy for hardware speedup.

Path-only conflict validation removes a previously hidden scale limit. The
explored-resource read set caused 1,206 of 13,458 \texttt{clma} connections to
overflow in the first round and still left 965 unrouted after 30 rounds.
Recording only candidate-path versions routes all 13,458 in round one and
reaches a legal solution in eight rounds. Host C-model routing time falls from
1,099.5\,s for the nonconvergent run to 154.7\,s for the legal run (7.11$\times$);
this is an algorithmic C-model comparison, not a U280 speedup claim.

Table~\ref{tab:pathreads-qor} reports the first four-case native-VPR matrix for
issue eight, every-round STA, and deterministic tail serialization after three
stagnant rounds. Every exported route passes VPR consistency and feasibility.
The geometric-mean wirelength and CPD deltas are +4.06\% and +1.48\%.

\begin{table}[t]
\centering
\small
\begin{tabular}{lrrrr}
\toprule
Circuit & Rounds & WL $\Delta$ & CPD $\Delta$ & Strict? \\
\midrule
\texttt{diffeq}  & 22 & +4.43\% & +0.58\% & yes \\
\texttt{frisc}   & 18 & +4.66\% & +0.88\% & yes \\
\texttt{elliptic}& 20 & +3.93\% & +3.86\% & no \\
\texttt{clma}    &  8 & +3.22\% & +0.66\% & yes \\
\bottomrule
\end{tabular}
\caption{Native VPR validation of the path-read-set policy. Strict denotes
WL$\leq$5\% and CPD$\leq$3\%; all four pass a relaxed 5\%/5\% throughput tier.}
\label{tab:pathreads-qor}
\end{table}

\subsection{Pre-registered board and ablation protocol}

To avoid selecting a favorable measurement after implementation, the final
board experiment is fixed as follows. Each case receives one warm-up and at
least five measured runs. CPU and U280 executions use identical packed netlist,
placement, RR graph, initial route, seed, channel width, timing constraints,
and termination limits. We report medians and dispersion for graph loading,
task preparation, Search, Commit, replay, VPR/Tatum update, export, and complete
wall time. Device power is integrated over the same interval; idle-subtracted
energy is reported per connection and per completed design.

The ablation matrix changes one mechanism at a time: (A0) full \system{}; (A1)
retransfer or reinitialize mutable state per iteration instead of persistence;
(A2) 1, 2, and 4 Search CUs; (A3) 1, 2, 4, and 8 queue banks; (A4) serialize
Search and Commit to remove speculation; (A5) remove replay by scheduling only
non-overlapping tasks; and (A6) timing updates every 1, 2, 4, or all iterations.
Every point must remain legal and report replay rate, wirelength, and CPD. The
primary acceptance target is at least 3$\times$ complete-routing-phase speedup over the multicore
reference, no more than 2\% mean wirelength or CPD degradation, and
bitwise-repeatable \system{} output. Missing these gates narrows the result to
an architecture study rather than the stated acceleration claim.

\subsection{Evidence boundary}

Table~\ref{tab:validation} establishes representation, control, RTL equivalence,
implementation, and legality. It does not establish the central quantitative
speedup target. A submission-ready accelerator evaluation must execute the
same full cases on a physical U280, including graph transfer, per-iteration
Search/Commit/replay, timing snapshots, and route export; repeat runs; record
energy; and ablate persistent state, CU count, MultiQueue banks, version checks,
and timing-update period. Those measurements are represented in the source as
an explicit unresolved requirement rather than estimated from hardware
emulation: \todoresult{end-to-end speedup, energy, and ablations}.

\subsection{Threats to validity}

The edge delay stored in the ABI is a local RC approximation and does not
exactly reproduce full route-tree Elmore coupling. VPR/Tatum still supplies the
real timing graph and constraints; native route re-analysis differs from the
feedback oracle by 0.011\,ns on \texttt{frisc} and 0.00434\,ns on
\texttt{elliptic}. Only clustered routing is supported by the
current oracle. The commercial FPGA24 case checks UltraScale+ legality, yet its
missing user timing constraints prevent a commercial timing claim. Finally,
the ten CPU pairs establish determinism for fixed seeds, not sensitivity across
seeds or physical-board run-to-run variance. The prescribed experiments
separate these questions.

\begin{table*}[t]
\centering
\small
\begin{tabular}{lccccc}
\toprule
System & Intra-search & Overlap & Deterministic & Cross-RRR state & Real STA \\
\midrule
Galois \cite{moctar2018galois} & yes & limited & yes & no & no \\
Connection-aware \cite{zhou2020connection} & no & partitioned & not central & no & yes \\
MultiQueue \cite{singer2024multiqueue} & yes & n/a & implementation-dependent & no & yes \\
CUFR \cite{zang2024cufr} & no & minimized & yes & CPU resident & weak \\
Potter \cite{zang2024potter} & no & yes & not central & CPU resident & yes \\
\system{} & yes & yes & fixed commit & HBM resident & VPR/Tatum \\
\bottomrule
\end{tabular}
\caption{Mechanism-level positioning. Persistent state here specifically means
accelerator state retained across complete RRR iterations.}
\label{tab:related}
\end{table*}

\section{Correctness and Performance Accounting}

\subsection{Serial equivalence of an epoch}

The transaction protocol provides a simple refinement argument. Let $S_i$ be
the authoritative resource and route-tree state before candidate $i$ in the
fixed Commit order. Search records the version of every mutable resource it
reads and proposes a path but does not change $S_i$. If all recorded versions
still match, Commit applies the candidate as one logical state transition,
yielding $S_{i+1}$. Otherwise, it leaves $S_i$ unchanged and places the task in
the replay set. Consequently, every accepted candidate observes the same
versions that a serial execution at its commit point would observe. Rejected
candidates have no visible partial update. Induction over the fixed order gives
one reproducible serial history for the accepted candidates in an epoch.

This property does not assert that speculative execution reproduces VPR's
exact path choices. The banked MultiQueue deliberately changes tie breaking,
and replay changes when a connection is searched. The required invariant is
instead semantic: each committed path is connected and in range, occupancy is
the sum of committed uses, each connection has one authoritative route-tree
representation, and termination requires zero overuse. The independent output
validator checks these invariants after every test. Deterministic ordering then
makes failures replayable and permits byte-level comparison between the
C-model and RTL.

Versions protect all resource fields whose changes can invalidate a candidate,
including occupancy-dependent cost and ownership. Immutable topology and delay
records require no versions. Rip-up and history-cost updates occur at epoch
boundaries, after Commit drains, so Search never observes a partially updated
negotiated-congestion round. Timing feedback follows the same boundary: route
changes are exported, VPR/Tatum analyzes the complete snapshot, and the next
epoch receives all connection criticalities together. Thus transactional
isolation covers the accelerator phase while an explicit barrier covers global
PathFinder and STA updates.

\subsection{What end-to-end speedup must include}

A kernel-only ratio would hide the costs that persistence is intended to
remove. For board run $r$, we therefore define accelerator time as
\begin{equation}
 T_{\mathrm{acc}}^{(r)}=T_{\mathrm{load}}+\sum_j
 (T_{\mathrm{search},j}+T_{\mathrm{commit},j}+T_{\mathrm{replay},j}
 +T_{\mathrm{STA},j})+T_{\mathrm{export}},
\end{equation}
where $j$ ranges over complete RRR iterations. Setup, buffer migration,
synchronization, and route validation remain in their corresponding terms.
The primary result will compare the median complete $T_{\mathrm{acc}}$ against
the median complete multicore VPR route phase; a second ratio will include
packing and placement. Search-kernel throughput is diagnostic only.
Balanced (3\% CPD/5\% wirelength) and relaxed Throughput (5\%/5\%) points form
separate Pareto tiers. Issue-eight/every-round-STA with path-only conflict
validation passes Balanced on three of four VTR cases and the relaxed tier on
all four; C-model time is not board-performance evidence, and neither tier
replaces the aspirational 2\% gate.

The CPU data in Table~\ref{tab:cpu} already show why both ratios are
needed. Sixteen workers improve the routing stage by 1.356$\times$ in geometric
mean, yet the whole flow improves by only 1.046$\times$. Persisting the RR graph
can remove repeated graph migration, but host STA and serial front-end work
remain. We will therefore report the fraction of time in each term above, PCIe
bytes per iteration, accepted candidates per Search wave, version-conflict
rate, mean replay count, and achieved HBM bandwidth. These counters distinguish
compute saturation from Commit serialization, poor locality, conflict-heavy
speculation, and host-side limits.

Energy uses the same boundary. Board power is sampled from before the initial
load through final export, synchronized with host timestamps. Idle power over
an equal interval is subtracted and both raw and corrected joules are retained.
Reporting joules per legal routed connection alongside total joules prevents a
faster but incomplete run from appearing efficient. Runs that time out, violate
QoR gates, fail validation, or omit a timing update are failures rather than
samples removed from the distribution.

\section{Related Work}

PathFinder \cite{mcmurchie1995pathfinder} and VPR
\cite{betz1997vpr,vtr8} define the negotiated-congestion academic baseline.
Deterministic Galois execution studies both inter-net and intra-search CPU
parallelism \cite{moctar2018galois}; connection-aware partitioning increases
granularity and improves Titan23 QoR/runtime \cite{zhou2020connection}. These
works prevent us from claiming determinism or connection-level parallelism
alone as novel.

MultiQueue routing demonstrates relaxed A* priority ordering for parallel
wave expansion \cite{singer2024multiqueue}. \system{} uses a hardware banked
MultiQueue but contributes the surrounding persistent state and transaction
protocol. For commercial devices, RWRoute provides an open timing-driven
UltraScale+ baseline \cite{zhou2021rwroute}. CUFR partitions nets with a
recursive ternary tree \cite{zang2024cufr}; Potter tolerates overlapping
partitions \cite{zang2024potter}; and AceRoute adaptively enables
intra-connection exploration \cite{wei2024aceroute}. Thus overlap tolerance
and two-level parallelism are also prior art. \system{}'s claim is their
hardware realization as versioned deterministic transactions over state that
persists across complete RRR and real-timing iterations.

Earlier FPGA-assisted routing work demonstrates that path search can benefit
from spatial pipelines or CPU--FPGA cooperation \cite{korolija2019fpga}. Such
kernel acceleration is complementary to our focus on state lifetime,
authoritative commit, and complete control-loop validation.

\section{Conclusion}

\system{} reorganizes timing-driven FPGA routing around persistent HBM state,
speculative Search CUs, and deterministic version-checked Commit. The
production ABI, real VPR/Tatum timing feedback, full U280 implementation, and
byte-identical multi-round RTL validation establish a coherent routing system
rather than an isolated maze kernel. Ten deterministic VPR pairs characterize
the multicore reference and expose the whole-flow bottleneck. Native VPR now
validates four path-read-set cases, including a 7.11$\times$ C-model scalability
recovery on \texttt{clma}. The remaining steps are physical-board end-to-end
time and energy, the six non-VTR QoR cases, and the registered ablations.

\clearpage
\bibliographystyle{ACM-Reference-Format}
\bibliography{references}

\end{document}
